\section*{Note to the Reader}
\addcontentsline{toc}{section}{Note to the Reader}

This is not the usual kind of ``computer'' book.

You won't find chapters on hash tables, compiler passes, or how to get a job
writing backend services. You also won't find clean separation between physics,
computer science, and ``AI safety.'' \COMPUTER{} is written on the assumption
that those separations are mostly accidents of history.

The working hypothesis here is blunt:

\begin{quote}
Computation is what reality is made of, not what we do \emph{to} reality.
\end{quote}

Everything else follows from taking that seriously.

\medskip

\noindent
\textbf{How to read this.}

\begin{itemize}
  \item \textbf{Parts I–III} build the ontology and the ``physics'' of \COMPUTER{}:
        Recursive Meta-Graphs (RMGs), double-pushout rewrite, MRMW, curvature,
        superposition as rewrite bundles, measurement as minimal path collapse.

  \item \textbf{Part IV} shows you machines that only make sense once you accept
        multiversal computation as the default: time-travel debugging,
        counterfactual execution engines, adversarial universes, and deterministic
        optimization across worlds.

  \item \textbf{Part V} sketches an architecture that could actually exist:
        a C$\Omega$MPILER and runtime capable of hosting those machines without
        lying about what they’re doing.

  \item \textbf{Part VI} is the jump: treating our physical universe as just
        another RMG (C$\Omega$SMOS), and intelligence and ethics as geometry problems
        in MRMW.
\end{itemize}

You can read linearly, or you can skim until something catches and then
backfill the definitions from the C$\Omega$DEX. The book tries to be self-similar:
concepts repeat in different guises; the same diagrams reappear at different scales.

\medskip

\noindent
\textbf{What this book is not.}

\begin{itemize}
  \item It is not a proof that ``the universe is a computer.'' It is a concrete
        model of what that claim would \emph{mean}, with enough structure that
        you can try to break it.

  \item It is not a complete theory. There are open conjectures and unproven
        bridges everywhere. That’s intentional.

  \item It is not neutral. It has opinions about how we should build future systems,
        what we should be terrified of, and which kinds of universes we should refuse
        to inhabit.
\end{itemize}

\medskip

If you are a physicist, expect parts of this to feel like an API spec for the
laws you already know. If you are a computer scientist, expect ``the machine''
to suddenly include galaxies. If you are an ML person, expect models to shrink
back down to what they are: ways of steering flows through a much larger space.

If you are none of those, but you are suspicious that reality is more like a
badly documented distributed system than a smooth manifold—this is your book.

\begin{flushright}
\COMPUTER{} is the name we’re giving to that suspicion,\\
and a proposal for what to build on top of it.
\end{flushright}
